% !TeX root = ../main.tex
\section{Generation 2: a clean null across three training oracles}
\label{sec:gen2}

Generation~2 kept the method and changed the apparatus: a 1B therapist instead of 7B, harder
patients, a stronger oracle with a JSON-schema rubric, and three training oracles swept at both
look-ahead depths for five matched iterations each. Because every model was scored on all six
questionnaires, it supplies $105$ persona-paired $\K$ contrasts --- the widest single test of
look-ahead available anywhere in this work.

The result is a null, and an unusually clean one.

\begin{table}[t]
\centering\small
\begin{tabular}{@{}lrrrrr@{}}
\toprule
\textbf{Training oracle} & \textbf{Iters} & $\K{=}5$ \textbf{ahead} & \textbf{mean} $\dlt$ & \textbf{mean} $\dz$ & \textbf{Holm sig.} \\
\midrule
Q1{+}Q2 & 5 & 4/5 & $-0.060$ & $-0.06$ & 0 \\
WAI-SR  & 5 & 2/5 & $+0.020$ & $+0.02$ & 0 \\
CSQ-8   & 5 & 3/5 & $-0.018$ & $-0.02$ & 0 \\
\midrule
\textbf{all three, \QQ{}} & 15 & \textbf{9/15} & $\mathbf{-0.019}$ & $-0.02$ & \textbf{0} \\
\textbf{all three, 7 metrics} & 105 & \textbf{61/105 (58.1\%)} & --- & --- & \textbf{0} \\
\bottomrule
\end{tabular}
\caption{Generation~2. $\dlt = \K_0 - \K_5$, persona-paired, $n = 96$ per cell; negative
favours look-ahead. Not one of the $105$ contrasts survives Holm correction in either
direction, and $\K{=}5$ wins about as often as a coin.}
\label{tab:gen2}
\end{table}

\subsection{Reading the null}
Three features make this more informative than a typical failure to reject.

\paragraph{The effect sizes are not merely non-significant, they are negligible.}
Mean $\dlt$ on the \QQ{} composite is $-0.019$ --- an order of magnitude smaller than
generation~1's $-0.132$, and below the oracle's own measurement noise (mean absolute
test--retest difference $0.04$--$0.09$). Per-metric means across the whole sweep range from
$-0.033$ to $-0.008$. This is not an underpowered version of generation~1's effect; it is the
absence of one. With $96$ paired observations per cell, a true effect the size of
generation~1's would have been detected comfortably.

\paragraph{It is a coin flip, not a trend.}
Look-ahead leads in $61$ of $105$ contrasts ($58.1\%$) and trails in the rest. Across the three
training oracles the sign is inconsistent: Q1{+}Q2 and CSQ-8 tilt slightly toward $\K{=}5$,
WAI-SR slightly away. A real effect attenuated by a harder setting would show a consistent sign
with shrunken magnitude; what generation~2 shows is no reliable sign at all.

\paragraph{It replicates across three independent reward definitions.}
The three training oracles are genuinely different objectives --- session satisfaction plus
working alliance, the Working Alliance Inventory, and the Client Satisfaction Questionnaire.
Look-ahead fails to help under all three. Whatever removed the effect is not specific to the
reward the ICLR experiment happened to optimize.

\subsection{What generation 2 does and does not establish}
It establishes that look-ahead's benefit is not a property of the PTO method as such: the same
method, with the same lever, produced a large effect in one setting and nothing in another.

It does not, on its own, tell us the effect has \emph{reversed} --- a null is compatible with a
small positive effect that this sweep could not resolve. That is what generation~3 settles.

It is also worth recording what generation~2's myopic arms did, because \S\ref{sec:discussion}
uses it: unlike generation~1, generation~2's $\K{=}0$ arms \emph{did} learn. Starting from an
untrained base of $2.378$ on \QQ{}, the myopic arms reached $2.770$ (Q1{+}Q2 oracle), $2.596$
(WAI-SR) and $2.629$ (CSQ-8) after five iterations --- gains of $+0.22$ to $+0.39$, against
generation~1's $+0.023$. The setting where look-ahead stopped helping is also the setting where
the one-step reward started working.
